% Fragmento público generado desde propietarios íntegros.
% Residencia: 52.7.2.
% Fuente: ../incoming/radion_monografia_20260827/manuscrito/sections/08_complejo_electromagnetismo.tex:59-158
\subsubsection{Constitutive Operator on the Channels \(90/120\)}

\paragraph{Common and Oriented Component Decomposition}

We define
\begin{align}
\mathcal E
&=\log\bigl[(1-q_+^{90})(1-q_-^{90})\bigr]
-\log\bigl[(1-q_+^{120})(1-q_-^{120})\bigr],
\label{eq:Ecommon}\\
\mathcal B
&=\log\frac{1-q_-^{90}}{1-q_+^{90}}
-\log\frac{1-q_-^{120}}{1-q_+^{120}}.
\label{eq:Borient}
\end{align}
Under \(C^*\mapsto-C^*\), \(q_+\leftrightarrow q_-\) are interchanged. Therefore,
\(\mathcal E\) is even and \(\mathcal B\) is odd.

The normalized constitutive equations are
\begin{equation}
\widehat\mu=e^{\mathcal E+\mathcal B},
\qquad
\widehat\varepsilon=e^{\mathcal E-\mathcal B},
\qquad
\widehat Z=e^{\mathcal B},
\qquad
\widehat c=e^{-\mathcal E}.
\label{eq:constitutives}
\end{equation}
They are verified exactly.
\begin{equation}
\widehat Z=\sqrt{\frac{\widehat\mu}{\widehat\varepsilon}},
\qquad
\widehat c=(\widehat\mu\widehat\varepsilon)^{-1/2}.
\label{eq:constitutive-identities}
\end{equation}
Conjugation produces
\begin{equation}
C^*\mapsto-C^*:\qquad
\widehat\mu\leftrightarrow\widehat\varepsilon,
\quad
\widehat Z\leftrightarrow\widehat Z^{-1},
\quad
\widehat c\mapsto\widehat c.
\label{eq:EM-conjugation}
\end{equation}

\begin{theorem}[Constitutive polarity]
The reader \(90/120\) converts the common mode and oriented contrast of the
pair \((A,C^*)\) into two conjugate constitutive quantities. Sheet inversion
exchanges permeability and permittivity, reverses impedance, and preserves normalized
velocity.
\end{theorem}

This is the precise formulation of the intuition \enquote{electrical part and
magnetic part}. One writes neither \(A=E\) nor \(C^*=B\). Instead, one writes the map
\[
(A,C^*)\to(x,y)\to(q_+,q_-)
\to(\mathcal E,\mathcal B)
\to(\widehat\varepsilon,\widehat\mu,\widehat Z,\widehat c).
\]

\paragraph{Elastic-rotational realization of the constitutive polarity}

The same information may be published in the complex regime by
\begin{equation}
Z_{\mathrm{em}}=e^{\mathcal E/2}
\left(\cos\frac{\mathcal B}{2}\,I
-\sin\frac{\mathcal B}{2}\,J\right).
\label{eq:Zem}
\end{equation}
The factor \(e^{\mathcal E/2}\) is an elastic dilation; the factor generated by \(J\) is a rotation. Complex notation does not add a mysterious imaginary dimension to a previously given real world: it records two operational modes already separated by the TRIT.

\paragraph{Transduction of Orientation into Effective Load}

The radion sheet \(\sigma\) fixes an orientation before a
unit of charge is introduced. Dimensional Mechanics then provides the transducer
\[
e_{\mathrm{int}}=\sqrt{\frac{2\alphah\hfull}{Z_{\mathrm{int}}}},
\qquad
Q=n_Qe_{\mathrm{int}}.
\]
The causal chain is
\[
\text{sheet}\to\text{orientation}\to\text{constitutive}
\to\text{dimensional transducer}\to\text{physical charge}.
\]
Calling the unit \emph{radion} does not prematurely identify the orientation
with the coulomb; it makes visible the locus where the charge sign may reside.

\begin{narrativebox}
Electricity and magnetism are not affixed to the circle from outside. They emerge
when the same state is read through product and quotient: the product
retains the common mode, while the quotient retains the orientation. The LC oscillator
realizes the exchange; the reader \(90/120\) publishes the constitutive relations; the
dimensional chart assigns units.
\end{narrativebox}
